% !TeX root = ../main.tex
\section{Introduction}
\subsection{Background and Motivation}
The use of advanced data analytics methods in sport gained popularity in the mid-1990s through the use of sabermetric principles in professional baseball \cite{Moneyball}. Since then, data-driven analysis has become central to performance evaluation, tactical decision-making, and modern sports broadcasting. During live coverage, high-level statistics such as possession, shot counts, and pass completion rates are commonly shown alongside the scoreline and elapsed time. Historically however, the collection of this data has typically been outsourced to specialised and expensive third-party analytics providers. For example, major English and Scottish football leagues have long-term data rights agreements with Opta, granting exclusive responsibility for the collection and distribution of official match statistics across broadcast, media, and analytical platforms \cite{StatsPerform}.

More recently, the growth of online streaming platforms has increased the broadcast coverage of lower-budget and regional sporting events \cite{pwc_sports_streaming_2025}. At the same time, amateur and semi-professional teams are increasingly adopting data-driven approaches to support training and performance evaluation \cite{futuremarketinsights_sports_analytics_2025}. Advances in computer vision now make it possible to analyse sports video streams in real time, enabling low-cost, on-site generation of match statistics without reliance on cloud infrastructure. This could reduce costs for both sports broadcasting and sports science applications, with particular benefits for lower-budget organisations.
\subsection{Current Approaches and Research Gap}
 Much of existing work on this topic has prioritised analytical accuracy and offline system design over deployment on constrained hardware. This results in many proposed systems relying on relatively large models, multi-stage processing, or GPU resources that are not realistic for compact edge devices, leaving a gap between strong research prototypes and systems that are practical for low-cost, real-time use in amateur clubs/broadcast settings, or field-side analysis. There is also limited work examining how optimisation, compression, and end-to-end system design interact in soccer analytics when latency, temperature, power, and GPU/CPU usage must be considered alongside model accuracy. These gaps in current approaches, which will be explored in greater detail in the literature review, motivate the problem aims and scope targeted by this thesis.

\subsection{Problem Aims}

These trends motivate the aim of this work: \textbf{the optimisation and compression analysis of embedded computer vision models for real-time soccer analysis on an edge device.} The specific objectives are as follows: 

\subsubsection{Model Training}
\label{sec:problem1}
 The first and most basic objective of this project is to train a number of computer-vision models of varying architecture and size, with the purpose of automatically extracting low-level information from video of soccer games. The model architecture and size should be selected with the intention of being run on an edge device, but should be trained to have an $\text{mAP}_{50-95}$ within 5\% to that of other modern, non real-time approaches \cite{tracking, rustebakke2025extracting, footandball}. Mean Average Precision (mAP) is a standard metric for localisation that evaluates performance by combining precision and recall. $\text{mAP}_{50\text{-}95}$ denotes the mean Average Precision calculated over IoU/OKS (Intersection over Union/Object Keypoint Similarity) thresholds from 0.50 to 0.95, providing a more comprehensive assessment of localisation accuracy than a single-threshold metric such as $\text{mAP}_{50}$.




\subsubsection{Model Compression}
\label{sec:problem2} The second objective is to effectively compress the trained models, aiming to reduce inference latency, CPU/GPU utilisation, device temperature and power draw without significantly affecting analytical performance. Specifically, the compressed models should perform to within 15\% $\text{mAP}_{50-95}$ of modern computer vision models used for soccer analysis \cite{tracking, rustebakke2025extracting, footandball}, while utilizing no more than 80\% of the chosen edge device's total CPU, and GPU availability. For this project, an NVIDIA Jetson Orin Nano Dev Kit will be used as the target edge device \cite{JetsonOrinNano}. Similarly, the power draw and temperature should not exceed 80\% of the device's recommended limits. NVIDIA recommends a maximum CPU and GPU temperature of $85^\circ$ and a max power draw of 25W, and so $80\%
\times85^\circ=68^\circ$ and $80\%
\times25W=20W$ will used as the initial recommended limits \cite{JetsonOrinNano}. 

\subsubsection{Full Stack System Proposal} The final objective is to propose a full-stack system design for real-time soccer analysis on an edge device, utilising the trained and compressed models. This will include pre- and post-inference processing to determine pitch locations, track players, assign teams, and extract high-level statistics such as possession, pass counts, and shot counts. The system should be implemented to minimise latency, with a minimum target of 100 ms (10 FPS) and an ideal target of 50 ms (20 FPS). It should generate high-level statistics, such as pass counts, within 15\% of the accuracy of modern computer vision models used for soccer analysis \cite{tracking}, while maintaining the same 80\% CPU/GPU utilisations, temperature limits, and power draw budget as the second objective.

